\section{Группы}

\begin{defn}{Множество с бинарной операцией}{}
    Множество \( M \), с заданным отображением \( M \times M \to M \) \( (a, b) \mapsto a \circ b \)

    Обозначается \( (M, \circ) \)
\end{defn}

\begin{defn}{Полугруппа}{}
    Множество с бинарной операцией \( (M, \circ) \) называется полугруппой,
    если эта операция ассоциативна:
    \[
        \forall a, b, c \in M \ (a \circ b) \circ c = a \circ (b \circ c)
    \]

    Обозначается \( (S, \circ) \)
\end{defn}

\begin{example}{Не всякая операция ассоциативна}{}
    \begin{itemize}
        \item
            \( M = \NN \) и \( a \circ b = a^b \) --- не ассоциативна:
            \[
                2^{(1^2)} \neq (2^1)^2
            \]
        \item
            \( M = \ZZ \) и \( a \circ b = a - b \) --- тоже не ассоциативна
    \end{itemize}
\end{example}

\begin{defn}{Моноид}{}
    Полугруппа \( (S, \circ) \) называется моноидом,
    если в ней существует нейтральный элемент:
    \[
        \exists e \in S: \ e \circ a = a \circ e = a
    \]
\end{defn}

\begin{example}{}{}
    Является ли \( (\NN, +) \) моноидом?
    Ответ зависит от географического положения,
    ведь в России не \( 0 \) не считают натуральным,
    а, например, во Франции, наоборот.
    Далее в курсе \( 0 \) не является натуральным.
\end{example}

\begin{remark}{}{}
    Если нейтральный элемент существует, то он единственный:
    \[
        e_1 = e_1 \circ e_2 = e_2
    \]
\end{remark}

\begin{defn}{Группа}{}
    Моноид \( (S, \circ) \) называется группой,
    если \( \forall a \in S \: \exists b \in S : a \circ b = b \circ a = e \)

    Обозначается \( (G, \circ) = (G, \cdot) = G \).

    Чаще всего используется мультипликативная запись: \\
    \( (a, b) \mapsto ab\), \( e = 1 \), обратный элемент обозначается \( a^{-1} \).
\end{defn}

\begin{exercise}{}{}
    Если обратный элемент существует, то он единственный.
\end{exercise}

\begin{defn}{Абелева группа}{}
    Группа \( G \) коммутативна (абелева), если \( \forall a, b \in G \ ab = ba \).

    В абелевых группах используется аддитивное обозначение: \\
    \( (a, b) \mapsto a + b\), \( e = 0 \), обратный элемент обозначается \( -a \).
\end{defn}

\begin{defn}{Порядок группы}{}
    Порядок группы \( G \) \( |G| \) --- число элементов в ней.

    \begin{itemize}
        \item
            \( |G| < \infty \Rightarrow G \) --- конечна
        \item
            \( |G| = \infty \Rightarrow G \) --- бесконечна
    \end{itemize}
\end{defn}

\begin{example}{Группы}{}
    \begin{enumerate}
        \item
            Числовые аддитивные:

            \( (\ZZ, +) \), \( (\QQ, +) \), \( (\RR, +) \), \( (\CC, +) \), \( (\ZZ_n, +) \)
        \item
            Числовые мультипликативные:

            \( (\QQ \setminus \{ 0 \}, \times) \),
            \( (\RR \setminus \{ 0 \}, \times) \),
            \( (\CC \setminus \{ 0 \}, \times) \),
            \( (\ZZ_p \setminus \{ 0 \}, \times) \)
        \item
            Группы матриц:

            \( GL_n (\RR) = \{ A \in M_n(\RR) \ \vline \ \det A \neq 0 \} \) --- общая линейная

            \( SL_n (\RR) = \{ A \in M_n(\RR) \ \vline \ \det A = 1 \} \) --- специальная линейная

            Есть еще куча разных групп: верхнетреугольные / нижнетреугольные / диагональные
            (конечно же, нужно взять только обратимые)
        \item
            Группы перестановок:

            \( S_n \), \( |S_n| = n! \)

            \( A_n \) --- четные перестановки, \( |A_n| = \frac{n!}{2} \)
    \end{enumerate}
\end{example}

\begin{exercise}{}{}
    \begin{itemize}
        \item
            \( S_n \) коммутативна \( \Leftrightarrow n \leq 2 \)
        \item
            \( A_n \) коммутативна \( \Leftrightarrow n \leq 3 \)
    \end{itemize}
\end{exercise}

\begin{defn}{Подгруппа}{}
    Подмножество \( H \) группы \( G \) называется подгруппой,
    если (определения эквивалентны):
    \begin{itemize}
        \item
            \begin{enumerate}
                \item
                    \( H \neq \varnothing \)
                \item
                    \( \forall a, b \in H \: ab^{-1} \in H \).
            \end{enumerate}
        \item
            \begin{enumerate}
                \item
                    \( e \in H\)
                \item
                    \( \forall a, b \in H \: ab \in H \)
                \item
                    \( \forall a \in H \: a^{-1} \in H \)
            \end{enumerate}
    \end{itemize}
\end{defn}

\begin{remark}{}{}
    В любой группе есть две подгруппы \( H = \{ e \} \) и \( H = G \)

    Они называются несобственными.
\end{remark}

\begin{prop}{}{}
    Любая подгруппа в \( (\ZZ, +) \) имеет вид \( k \ZZ \) (\( k \in \ZZ_{\geq 0} \))
\end{prop}

Очевидно, что \( k \ZZ \) --- подгруппа.

Пусть \( H \subseteq \ZZ \) --- подгруппа.

Если \( H = \{ 0 \} \), то \( k = 0 \).

Иначе в \( H \) есть ненулевой элемент.
Тогда в \( H \) есть положительный элемент.
Пусть \( k \) --- наименьший положительный элемент в \( H \).
Тогда \( k \ZZ \subseteq H \).
Пусть \( a \in H \).
Поделим с остатком: \( a = kq + r \), где \( 0 \leq r < k \).
\( r = a - kq \). \( a \in H, k \in H \Rightarrow r \in H \).
\( k \) наименьшее положительное число в \( H \), значит \( r = 0 \),
то есть \( a \in k \ZZ \).

\begin{defn}{Циклическая подгруппа}{}
    Пусть \( G \) --- группа, \( g \in G \).

    Циклической подгруппой,
    порожденной элементом \( g \) называется
    \[
        H = \langle g \rangle = \{ g^m, m \in \ZZ \}
    \]
\end{defn}

\begin{defn}{Порядок элемента}{}
    Пусть \( G \) --- группа, \( g \in G \)

    Порядком элемента \( g \) называется наименьшее натуральное число \( m \),
    для которого \( g^m = e \) или бесконечность, если такого \( m \) не существует.

    Обознается \( \ord (g) \)
\end{defn}

\begin{prop}{}{}
    \[
        | \langle g \rangle | = \ord (g)
    \]
\end{prop}

Пусть \( g^n = g^m \) и \( n > m \Rightarrow g^{n - m} = e \Rightarrow \) порядок \( g \) конечен.

Отсюда можно сделать вывод, что если \( \ord (g) = \infty \), то \( g^n \neq g^m \) при \( n \neq m \),
а значит и \( | \langle g \rangle | = \infty \).

Теперь пусть \( \ord (g) = m \).
По определению \( g^m = e \) и \( g^k \neq e \) при \( 0 < k < m \).
Поэтому \( g^0, \ldots, g^{m - 1} \) попарно различны (иначе можно применить трюк выше).
Любое целое число можно поделить с остатком на \( m \): \( s = mq + r \), где \( 0 \leq r < m \).
Но тогда понятно, что \( g^s = g^r \).
А значит \( \langle g \rangle = \{ g^0, g^1, \ldots, g^{m - 1} \} \Rightarrow |G| = m \).

Что и требовалось доказать.

\begin{defn}{Циклическая группа}{}
    Группа \( G \) назыается циклической, если \( \exists g \in G : G = \langle G \rangle \)
\end{defn}

\begin{remark}{}{}
    Нетрудно убедиться, что тогда \( G \) коммутативна и не более чем счетна.
\end{remark}

\begin{example}{}{}
    \begin{itemize}
        \item
            \( (\ZZ, +) = \langle 1 \rangle = \langle -1 \rangle \)
        \item
            \( (\ZZ_n, +) = \langle \ol{1} \rangle = \langle \ol{k} \rangle \), где \( (k, n) = 1 \)
        \item
            \( (\CC \setminus \{ 0 \}, \times) \supset \langle \varepsilon \rangle \),
            где \( \varepsilon \) --- первообразный корень
    \end{itemize}
\end{example}

\begin{defn}{Левый смежный класс}{}
    Пусть \( G \) --- группа, \( H \subseteq G \) --- подгруппа.

    Левый смежный класс элемента \( g \in G \) по подгруппе \( H \) --- это подмножество,
    которое обозначается \( gH \) и равно \( \{ gh \: \vline \: h \in H \} \).
\end{defn}

\newpage

\begin{lemma}{}{}
    Либо \( g_1 H = g_2 H \), либо \( g_1 H \cap g_2 H = \varnothing \).
\end{lemma}

Если их пересечение пусто, то условие верно, иначе есть общий элемент:
\[
    g_1 h_1 = g_2 h_2 \Rightarrow g_1 = g_2 h_2 h_1^{-1}
\]

\( H \) --- подгруппа, поэтому \( h_2 h_1^{-1} \in H \Rightarrow g_1 \in g_2 H \Rightarrow g_1 H \subseteq g_2 H \).

Аналогично \( g_2 H \subseteq g_1 H \), откуда \( g_1 H = g_2 H \).

\begin{lemma}{}{}
    \[
        |gH| = |H|
    \]
\end{lemma}

\( gH = \{ gh \: \vline \: h \in H \} \),
поэтому давайте сопоставим \( h \in H \) элемент \( gh \).
Откуда \( |H| \geq |gH| \).

Осталось показать, что никакие два элемента в \( gH \) не совпадают:
\[
    g h_1 = g h_2 \Rightarrow h_1 = h_2
\]

А значит \( |H| = |gH| \), что и требовалось доказать.

\begin{defn}{Индекс подгруппы}{}
    Индекс подгруппы \( H \) в группе \( G \) --- это
    \[
        [ G : H ] = | G / H | = \text{число левых смежных классов}
    \]
\end{defn}

\begin{thrm}{}{}
    Если \( G \) --- конечная группа, а \( H \) в ней подгруппа,
    то \( |G| = |H| \cdot [G : H] \)
\end{thrm}

Явно следует из двух предыдущих лемм.

\begin{coroll}{}{}
    Если \( G \) конечно, и \( H \subseteq G \) --- подгруппа,
    то \( |G| \vdots |H| \).
\end{coroll}

\begin{coroll}{}{}
    Если \( G \) конечно, и \( H \subseteq G \) --- подгруппа,
    то \( \forall g \in G \: |G| \vdots \ord(g) \).
\end{coroll}

Явно следует из \( \ord (g) = | \langle g \rangle | \)

\begin{coroll}{}{}
    \[
        g^{|G|} = e
    \]
\end{coroll}

Явно следует из того, что \( |G| \vdots \ord (g) \)

\begin{coroll}{}{}
    Малая теорема Ферма
\end{coroll}

\begin{coroll}{}{}
    Если \( |G| = p \) --- простое, то \( G \) --- циклическая группа,
    порожденная своим любым неединичным элементом.
\end{coroll}

Пусть \( g \in G \) и \( g \neq e \).

Тогда \( p = |G| \vdots | \langle g \rangle | \).
Так как \( | \langle g \rangle | \geq 2 \), \( | \langle g \rangle | = p = |G| \),
а значит \( |G| = \langle g \rangle \).

\begin{defn}{Правый смежный класс}{}
    Пусть \( G \) --- группа, \( H \subseteq G \) --- подгруппа.

    Правый смежный класс элемента \( g \in G \) по подгруппе \( H \) --- это подмножество,
    которое обозначается \( Hg \) и равно \( \{ hg \: \vline \: h \in H \} \).
\end{defn}

\begin{remark}{}{}
    Для них верны те же свойства, что и для левых смежных классов.
\end{remark}

\begin{defn}{Нормальная подгруппа}{}
    Подгруппа \( H \) нормальна, если \( gH = Hg \) \( \forall g \in G \).
\end{defn}

\begin{prop}{}{}
    Для подгруппы \( H \subseteq G \) следующие условия эквивалентны:
    \begin{enumerate}
        \item
            \( H \) нормальна
        \item
            \( gHg^{-1} \subseteq H \ \forall g \in G \)
        \item
            \( gHg^{-1} = H \ \forall g \in G \)
    \end{enumerate}
\end{prop}

\begin{itemize}
    \item
        \( 1 \Rightarrow 2 \)

        Рассмотрим \( h \in H, g \in G \)

        \( gH = Hg \Rightarrow gh = h'g \), \( h' \in h \)

        Тогда \( g h g^{-1} = h' g g^{-1} = h' \in H \)
    \item
        \( 2 \Rightarrow 3 \)

        \( g H g^{-1} \subseteq H \)

        \( \forall h \in H \) имеем \( h = (g g^{-1}) h (g g^{-1}) = g (g^{-1} h g) g^{-1} \in g H g^{-1} \)

        А это значит, что \( H \subseteq g H g^{-1} \Rightarrow H = g H g^{-1} \)
    \item
        \( 3 \Rightarrow 1 \)

        \( g H g^{-1} = H \)

        Домножим справа на \( g \): \( gH = Hg \)
\end{itemize}

\textbf{Цель}: определить факторгруппу, то есть ввести операцию на множестве левых смежных классов \( G / H \),
где \( G \) --- группа, а \( H \subseteq G \) --- подгруппа.

Хочется чего-нибудь интуитивного: \( (g_1 H) (g_2 H) = g_1 g_2 H \).

Проверим необходимые свойства:
\begin{enumerate}
    \item
        Ассоциативность: \( ( (g_1 H) (g_2 H) ) (g_3 H) = g_1 g_2 g_3 H = (g_1 H) ( (g_2 H) (g_3 H) ) \)
    \item
        Нейтральный элемент: \( eH = H \quad (eH) (gH) = (gH) = (gH) (eH) \)
    \item
        Обратный элемент: \( (gH)^{-1} = g^{-1} H \quad (g^{-1} H) (g H) = eH = (g H) (g^{-1} H) \)
\end{enumerate}

Осталось самое нетривиальное --- корректность операции.
Надо как-то убедиться, что операция не изменяется при выборе разных представителей \( g_1, g_2 \) смежных классов:
\begin{gather*}
    (g_1 h_2) (g_2 h_2) H = g_1 g_2 H
    \\
    h_1 g_2 (h_2 H) = g_2 H
    \\
    g_2^{-1} h_1 g_2 H = H
    \\
    \Updownarrow
    \\
    g_2^{-1} h_1 g_2 \in H
    \\
    \Updownarrow
    \\
    \text{\( H \) нормальна}
\end{gather*}

\begin{defn}{Факторгруппа}{}
    Факторгруппа группы \( G \) по нормальной подгруппе \( H \)
    называется множество (левых) смежных классов \( G / H \) с операцией
    \( (g_1 H) (g_2 H) = g_1 g_2 H \)
\end{defn}

\newpage

\begin{example}{}{}
    \begin{itemize}
        \item
            \( G = (\ZZ, +) \), \( H = n \ZZ \)

            \( G / H = \{ s + H = s + n \ZZ = \ol{s} \} \)

            \( G / H = \ZZ_n \)
    \end{itemize}
\end{example}
